% Derivation Performance Notes
% Source: docs/derivation_performance.md (migrated to LaTeX)
% Uses macros from preamble.tex

\section{Derivation Performance Notes}
\label{sec:derivation-performance}

\subsection{Spherical 4D Einstein-Maxwell + TT Gauge (theory\_radial.toml)}
\label{sec:spherical-4d}

Total wall time: ${\sim}40$ minutes. Peak memory: ${\sim}840$\,MB.

\subsubsection{Phase Breakdown}

\begin{table}[htbp]
\centering
\caption{Derivation phase breakdown for spherical 4D Einstein-Maxwell + TT gauge.}
\label{tab:spherical-phases}
\begin{tabular}{@{}llll@{}}
\toprule
Phase & Time (est.) & Memory delta & Notes \\
\midrule
xPert $\lag^{(2)}$ expansion & ${\sim}3$\,min & +5\,MB & Fast \\
$h$ decomp (6 components) & \textbf{${\sim}25$\,min} & +690\,MB & \textbf{DOMINANT} \\
\quad BatchedTBD per component & ${\sim}2$\,min & +3--5\,MB & 1000--2300 terms each \\
\quad Christoffel (first) & ${\sim}30$\,s & +285\,MB & Cache miss; subsequent ${\sim}5$\,s \\
$a$ decomp (4 components) & ${\sim}10$\,min & +50\,MB & Rank-1, smaller \\
ExportJSON & ${\sim}1$\,min & minimal & \\
\bottomrule
\end{tabular}
\end{table}

\subsubsection{Bottlenecks}

\begin{enumerate}
  \item \textbf{BatchedTraceBasisDummyWithMetric} (${\sim}20$\,min total for 10 components)
  \begin{itemize}
    \item Fused \texttt{TraceBasisDummy} + \texttt{Expand} + \texttt{MetricEval} in adaptive batch groups
    \item Generates 1000--2300 intermediate terms per component
    \item Single largest time sink---see ``Adaptive batch sizing'' below (\cref{sec:adaptive-batch})
  \end{itemize}

  \item \textbf{Memory growth} (154$\to$840\,MB during $h$ decomp)
  \begin{itemize}
    \item \texttt{Share[]} between components helps but doesn't fully reclaim
  \end{itemize}

  \item \textbf{Christoffel first-call cost} (+285\,MB, ${\sim}30$\,s)
  \begin{itemize}
    \item Already cached via \texttt{GetCachedChristoffels}---subsequent calls $O(1)$
    \item The 285\,MB is the Christoffel array itself ($4 \times 4 \times 4$ with symbolic entries)
  \end{itemize}
\end{enumerate}

\subsection{Flat 4D Einstein-Maxwell (theory.toml)}
\label{sec:flat-4d}

Total: ${\sim}40$ minutes. Peak: ${\sim}170$\,MB (Wolfram heap) / ${\sim}1.3$\,GB (OS RSS).

\begin{table}[htbp]
\centering
\caption{Derivation phase breakdown for flat 4D Einstein-Maxwell.}
\label{tab:flat-phases}
\begin{tabular}{@{}lll@{}}
\toprule
Phase & Time & Notes \\
\midrule
EOM $h$ decomposition (6 comp) & ${\sim}14$\,min & 332 \& 226 terms per comp \\
EOM $a$ decomposition (4 comp) & ${\sim}1.5$\,min & 99 terms per comp \\
Linearization total & ${\sim}16$\,min & Includes xPert + gauge \\
\textbf{Canonical Lagrangian decomp} & \textbf{${\sim}24$\,min} & \textbf{34 additive terms, DOMINANT} \\
Legendre + parse & $<1$\,s & Trivial \\
\bottomrule
\end{tabular}
\end{table}

The canonical Lagrangian decomposition is the dominant cost (24 of 40\,min). This is intrinsic to the full Legendre transform---each of the 34 Lagrangian terms must be independently decomposed via \texttt{DecomposeScalarExpression}. Previously (${\sim}10$\,min) the canonical pipeline was skipped; now it is always computed for correct coordinate-invariant energy measurement.

\subsection{Post-v6 Reference Timings (v0.33.9)}
\label{sec:post-v6-timings}

With the v6 perturbative-reduction iterative architecture (Parker--Simon, see \S\ref{sec:pr-v6}) and the accumulated optimisations below, representative derivation wall-times on a single-host devcontainer (single Wolfram license, cold shell, warm Mathematica cache) are:

\begin{table}[htbp]
\centering
\caption{Representative post-v6 derivation wall-times measured with \texttt{time uv run tidal derive} on a single-host devcontainer. Numbers are from v0.33.9; run-to-run variance is $\pm 10\%$ dominated by Mathematica kernel initialisation.}
\label{tab:post-v6-timings}
\begin{tabular}{@{}lllll@{}}
\toprule
Theory & Wall-time & $n_{\mathrm{dyn}}$ & Components & Notes \\
\midrule
\texttt{coupled\_scalars}       & ${\sim}14$\,s   & 2 & 2  & 1+1D, minimal cross-coupling \\
\texttt{euler\_heisenberg}      & ${\sim}19$\,s   & 1 & 4  & matter-only, $F^4$ correction (\#271) \\
\texttt{scalar\_potential\_well}& ${\sim}10$\,s   & 1 & 1  & fast-path precompute skip \\
\texttt{gertsenshtein}          & ${\sim}46$--$49$\,s & 2 & 6 & Einstein--Maxwell baseline \\
\texttt{graviton\_torsion}      & ${\sim}6$--$8$\,min & 2 & 13 & PGT, \texttt{--timeout 0} \\
\texttt{torsion\_gertsenshtein} & ${\sim}20$\,min & 3 & 20 & full PGT + EM + $\tilde{R}^2$ \\
\bottomrule
\end{tabular}
\end{table}

\paragraph{\texttt{euler\_heisenberg} commentary.} Before v0.33.9 this example could not derive at all --- first the \texttt{(F$\cdot$F)}$^2$ Power auto-distribution raised \verb|Validate::repeated + Throw[Null]|, and even with that fixed the CD ComponentValue precompute was skipped (the $\texttt{len(dyn\_fields)} \ge 2$ heuristic), leaving \verb|CD1eqca| shorthand unresolved in \verb|lagComp| so Component E--L found zero field functions. Both are fixed in commits 9d9e73f and 830a442; see \S\ref{sec:pr-eh-cd}. The $\approx 19$~s wall-time listed above is the first recorded measurement for this theory.

\paragraph{\texttt{gertsenshtein} commentary.} The v0.33.9 safety-net \verb|//.| of \verb|$CDShorthandReverseRules| before Component-E--L field detection was initially a $+24\%$ regression (45.6~s $\rightarrow$ 56.6~s). Gating it on \verb|!FreeQ[lagComp, _Symbol?(StringMatchQ["CD" ~~ DigitCharacter ~~ __]&)]| --- a no-op whenever the precompute already resolved everything --- restored the timing to $\approx 49$~s, within run-to-run variance of the 46~s baseline.

\subsection{Wolfram License Constraint}
\label{sec:license-constraint}

Only ONE \texttt{wolframscript} session at a time (engine license). Running two concurrent derivations exhausts the license and causes ``license error'' exit. Always serialize derivation commands. The license error can also occur mid-derivation if the WLS runs for too long; the Python post-processing (plane-wave reduction, constraint elimination) then fails because it checks \texttt{ret == 0}.

\textbf{Workaround}: If the WLS completes but Python post-processing is skipped due to license error, run the plane-wave reduction manually:

\begin{lstlisting}[style=python]
from tidal.symbolic.reduction import reduce_spec
import json, tomllib

spec = json.loads(Path('examples/data/file.json').read_text())
with open('examples/.../theory.toml', 'rb') as f:
    config = tomllib.load(f)
reduced = reduce_spec(spec, config['reduction'])
Path('examples/data/file.json').write_text(json.dumps(reduced, indent='\t'))
\end{lstlisting}

\textbf{Note}: Constraint elimination is now handled entirely by the Wolfram pipeline (before JSON export). There is no Python post-processing step for constraints.


\subsection{Optimizations}
\label{sec:optimizations}

\subsubsection{Timing Instrumentation}

All major derivation phases print \texttt{[TIMING]} lines to stderr with elapsed seconds. Phases instrumented:

\begin{itemize}
  \item \texttt{Linearization (xPert L\^{}(2) + EOM decomposition)}---overall linearization
  \item \texttt{EOM decomposition (<field>)}---per-field \texttt{DecomposeToComponents}
  \item \texttt{Canonical Lagrangian decomposition}---\texttt{DecomposeScalarExpression} loop
  \item \texttt{Legendre transform (momenta + H)}---canonical momentum + $H = \sum \pi v - L$
  \item \texttt{ParseHamiltonianExpression}---structured quadratic term parsing
\end{itemize}

Each canonical Lagrangian term also reports individual time and memory.

\subsubsection{Derivation Caching}

\texttt{tidal derive} hashes the generated WLS script (SHA-256) and stores it in the output JSON's \texttt{metadata.derivation\_hash}. On subsequent runs, if the TOML config produces an identical WLS script and the output JSON exists with a matching hash, \texttt{wolframscript} is skipped entirely.

\begin{lstlisting}[style=bash]
# First run: full derivation (~40 min for theory_radial.toml)
uv run tidal derive examples/gertsenshtein/theory_radial.toml

# Second run: cache hit, instant
uv run tidal derive examples/gertsenshtein/theory_radial.toml
# -> "Derivation cache hit: gertsenshtein_radial.json"

# Force re-derivation
uv run tidal derive examples/gertsenshtein/theory_radial.toml --force-derive
\end{lstlisting}

The hash captures ALL inputs: TOML fields, parameters, Wolfram pipeline code paths, gauge config, linearization options, reduction settings. Any change to the TOML or to the Python code that generates the WLS script invalidates the cache automatically.

\subsubsection{Adaptive Batch Sizing in BatchedTraceBasisDummyWithMetric}
\label{sec:adaptive-batch}

\texttt{BatchedTraceBasisDummyWithMetric} (\texttt{ComponentDecompose.wl}) fuses \texttt{TraceBasisDummy} + \texttt{Expand} + metric evaluation in batches to bound peak memory. Previously used a fixed batch size of 50.

Now uses \textbf{adaptive sizing}: the first batch runs at the default size (50), then the expansion factor (peak intermediate terms / input terms) is measured. If the expansion is high enough, subsequent batches are \textbf{decreased} in size to keep peak intermediate terms below ${\sim}2000$:

\begin{equation}
  \text{batch\_size} = \min\!\bigl(\text{default},\; \max(5,\; \lfloor 2000 / \text{expansion\_factor} \rfloor)\bigr)
  \label{eq:adaptive-batch}
\end{equation}

\textbf{Critical}: Batch size is NEVER increased above the default. \texttt{Expand[]} scales super-linearly with batch size due to cross-term interactions during simplification---empirically, doubling batch size can increase processing time by 10--20$\times$, not $2\times$. The default batch size of 50 is near-optimal for most theories.

\begin{itemize}
  \item Simple scalar theories (expansion ${\sim}1$--$5\times$): batch stays at 50 (default)
  \item Complex tensor theories (expansion ${\sim}40$--$50\times$): batches decrease to 20--30
  \item Per-batch diagnostics: \texttt{TraceBD}, \texttt{Expand}, \texttt{Eval} term counts + timing
\end{itemize}

Ref: \texttt{ComponentDecompose.wl}, \texttt{BatchedTraceBasisDummyWithMetric} function.

\subsubsection{EOM vs Canonical Decomposition (No Caching Possible)}

Investigation confirmed that the EOM pass decomposes the \textbf{equations of motion} (\texttt{VarD[L]}) via \texttt{DecomposeToComponents}, while the canonical path decomposes the \textbf{Lagrangian} itself via \texttt{DecomposeScalarExpression}. These are fundamentally different algebraic objects---the Legendre transform requires $L_\text{comp}$, not $\text{EOM}_\text{comp}$. The Wolfram kernel state (Christoffels, metric DownValues, background field DownValues) IS already reused between passes via \texttt{Share[]}.

\subsection{Component-Level E--L Performance (v0.15.0+)}
\label{sec:component-el-perf}

The component-level Euler--Lagrange pipeline (see \S\ref{sec:component-el} in Architecture) eliminates the \texttt{VarD} + \texttt{TraceBasisDummy} bottleneck for theories with many contracted abstract indices.

\begin{table}[htbp]
\centering
\caption{Derivation timing comparison for PGT torsion ($\torsion^2 + \tilde{R}^2$), 2+1D.}
\label{tab:component-el-timing}
\begin{tabular}{@{}llll@{}}
\toprule
Phase & Abstract VarD & Component E--L & Speedup \\
\midrule
xPert linearisation & ${\sim}200$\,s & ${\sim}60$\,s (with CD shorthand) & 3$\times$ \\
Phase A (component decomp) & --- & ${\sim}4$\,s & --- \\
Phase B (E--L per field) & $>$77\,min & $<$1\,s/field ($\sim$5\,s total) & $>$900$\times$ \\
\midrule
\textbf{Total} & $>$80\,min & ${\sim}65$\,s & $>$70$\times$ \\
\bottomrule
\end{tabular}
\end{table}

\textbf{Key insight:} The abstract VarD path scales with the number of \emph{contracted abstract indices} (for $\tilde{R}^2$: $\sim$45 indices, giving $O(\text{dim}^{n})$ scaling in \texttt{TraceBasisDummy}). The component E--L path operates on scalar expressions with no abstract indices, making it insensitive to the complexity of the original tensor structure.

The CD shorthand optimisation (pre-computing covariant derivative expansions) further reduces xPert linearisation from ${\sim}200$\,s to ${\sim}60$\,s, and is now applied to ALL theories (not just torsion).

\textbf{Note:} Component-level E--L is the default for all theories since v0.15.0. The abstract VarD path is still available as a fallback but is no longer used in the standard pipeline.

\subsection{Canonical Decomposition Optimisations (v0.18.6+)}
\label{sec:canonical-decomp-perf}

Two further optimisations reduce canonical Lagrangian decomposition time:

\paragraph{Single-term processing (v0.18.6).}
The previous approach split the second-order Lagrangian $\mathcal{L}^{(2)}$ into individual additive terms and processed each through \texttt{DecomposeScalarExpression} separately. The per-call overhead of $\sim$34 separate invocations exceeded the single-pass cost. Processing the full expression as one term reduces Gertsenshtein canonical decomposition from 66\,s to 43\,s (35\% faster). Torsion is unchanged ($\sim$9\,s) because \texttt{ExpandScalarWrappers} already pre-resolves Scalar wrappers individually.

\textbf{Bottleneck identified:} \texttt{TraceBasisDummy} dominates at 42.7\,s of 43\,s total, with $O(\text{dim}^{2K})$ scaling for $K \sim 4$ contracted index pairs per Einstein--Hilbert Ricci scalar product term.

\paragraph{Multi-field pipeline optimisations (v0.22.8).}
For combined multi-field theories (e.g.\ PGT + Einstein--Maxwell with 23 components), four extensible optimisations address a 15-minute derivation bottleneck:
\begin{enumerate}
  \item \textbf{Progressive $\mathcal{L}^{(2)}$ cleanup:} Interleaved \texttt{Expand} after each background-zeroing step, with adaptive per-term batched canonicalisation.
  \item \textbf{CD shorthand ComponentValue pre-computation:} Compute component arrays for CD1--CD4 shorthand tensors via \texttt{StaggeredToBasis}, reducing \texttt{TraceBasisDummy} dummy pairs from $K \sim 4$ to $\sim$0.
  \item \textbf{Continuous adaptive batching:} \texttt{DecomposeScalarExpression} batches large expressions with per-batch adaptive sizing, preventing blowup when later batches have different coupling structure.
  \item \textbf{Sector-aware splitting:} Automatically classify expanded Lagrangian terms by field content (graviton-only, torsion-only, EM-only, cross-sector), processing each sector independently.
\end{enumerate}

\subsection{Hamiltonian Phase Optimisations (v0.25.7+)}
\label{sec:hamiltonian-phase-perf}

For torsion theories with non-minimal couplings, the canonical Hamiltonian computation is the dominant bottleneck (Phase~B: IBP + Legendre transform). Two optimisations target Phase~B:

\paragraph{Pre-IBP torsion filter (Phase~B).}
Before the time-variable integration-by-parts loop, terms containing \emph{only} torsion component functions (no GW or EM field content) are removed from \texttt{lagComp}. These terms contribute zero to both canonical momenta ($\pi_h$, $\pi_a$) and the filtered Legendre subtraction ($-L$), so their removal is exact. For the \texttt{torsion\_gertsenshtein} theory, this eliminates ${\sim}40$--60\% of IBP work.

Ref: inspired by irreducible sector decomposition in HiGGS~\cite{Barker:2022kdk} (github.com/wevbarker/HiGGS), which computes constraints only for physically relevant torsion sectors. See also Hamilcar (github.com/wevbarker/Hamilcar) for index-free canonical analysis.

\paragraph{Collect before Legendre transform (Phase~B).}
\texttt{Collect[lagComp, fieldFuncList, Identity]} combines terms with identical field-function structure before the momentum computation loop. This reduces the number of \texttt{D[lagComp, vel]} evaluations by ${\sim}15$--25\%.

\paragraph{Phase~A caching (under investigation, \#228).}
Caching \texttt{DecomposeScalarExpression} results on abstract tensor expressions was tested but caused severe regression (${\sim}131$\,s $\to$ $>600$\,s timeout) due to expensive \texttt{Hash}/\texttt{SeparateScalarCoefficient} calls on deeply nested xAct expressions. Alternative caching strategies targeting post-decomposition (component scalar) expressions are being investigated; see \#228 for details.
